\section*{Open Questions}

The following five questions are genuinely open after this replication.
Each names a specific concrete probe that would close it.

\begin{enumerate}
  \item \textbf{Constructive composition with hardware error mitigation.}
        Does algorithmic Richardson-in-$1/N$ compose additively with ZNE
        or PEC, or do the two mitigations interfere in variance?
        The paper's Fig.~3(b,c) asserts composition works for one Pauli
        channel; we did not rerun those noise curves.
        \emph{Probe:} density-matrix Qiskit-Aer simulation of the 5-qubit
        TFIM with the paper's Pauli channel, sweeping
        $N\in\{15,20,25,30\}$ and gate-noise scaling $\{1,3,5\}\times$,
        comparing raw / algorithmic-only / physical-only / combined
        estimators for bias and variance.

  \item \textbf{Fault-tolerant relevance.}
        Is Richardson-in-$1/N$ a NISQ-era trick or does it also save
        $T$-gates in the fault-tolerant regime, where each ``shot'' is
        a full logical circuit and the $N$-fold classical repetition
        becomes a real cost multiplier?
        \emph{Probe:} $T$-count model for a 100-qubit chemistry Trotter
        circuit at $10^{-10}$ logical error; compare single high-$N$
        Trotter against 3-pt Richardson at smaller $N$s, both including
        the classical-repetition overhead, for matched end-to-end
        accuracy.

  \item \textbf{Interaction with product-formula alternatives (qDRIFT,
        qSWIFT).}
        Randomized product formulas convert algorithmic error into
        variance rather than bias. Can Richardson be layered on top,
        and does the combination beat either alone?
        \emph{Probe:} qDRIFT on the same 5-qubit TFIM; then a
        Richardson-on-qDRIFT variant extrapolating over an effective
        step parameter; report shots-to-$10^{-4}$-accuracy against
        first-order Trotter + Richardson at equal circuit depth.

  \item \textbf{Sensitivity to Trotter ordering.}
        The extrapolation coefficients $a_1,a_2$ depend on nested
        commutators whose numerical values are ordering-dependent
        (Lie vs.\ Strang, X-first vs.\ ZZ-first). Neither the paper
        nor this replication ablates ordering.
        \emph{Probe:} rerun 3-pt Richardson at $(15,20,25)$ for four
        orderings (Lie X-ZZ, Lie ZZ-X, Strang X-ZZ-X, Strang ZZ-X-ZZ);
        report the spread of mitigated error. Prediction: Strang
        orderings zero $a_1$ by symmetry, so Richardson becomes purely
        an $a_2$ removal.

  \item \textbf{Time-dependent Hamiltonians.}
        For $H(t)$, algorithmic error follows a Magnus expansion, not
        BCH, and may contain $\log(N)$ corrections that break the
        polynomial-in-$1/N$ assumption underlying Richardson.
        \emph{Probe:} 5-qubit TFIM with $B(t)=2+0.5\sin(\pi t/T)$,
        exact reference by fine Runge--Kutta on the 32-dim state;
        sweep $N\in\{10,15,20,25,30\}$ and test whether adding a
        $\log(N)/N$ term to the fit improves $\chi^2$
        significantly. If Richardson breaks, that is a real scope
        boundary the paper does not flag.
\end{enumerate}
